\section{Output: the flower JSON}\label{sec:output}
\stagebox{\statusbuilt}{\src{core/src/export.rs}, \src{core/src/bloom/export.rs}, \src{core/src/channels.rs}, \src{core/src/bin/automark.rs}, \src{viewer/FORMAT.md}, \src{viewer/src/mesh.ts}}

The core's one canonical output is a JSON document (format
\texttt{automark-flower}, version 1) shared with the viewer. The core owns the
geometry; the viewer only meshes, shades, animates and exports glTF.

\subsection{Contents}
\begin{itemize}[leftmargin=*]
\item \texttt{seed} as a decimal string (JavaScript numbers lose precision
  above $2^{53}$), \texttt{flower\_index}, units of the core radius and $+z$
  up.
\item \texttt{shells}: \texttt{aspect} ($\alpha$), \texttt{egg} ($\eta_e$),
  \texttt{point} ($\eta_p$), \texttt{slot} ($w$), \texttt{radius} (the layer
  radii $a_k$), \texttt{core} ($R_0$) and \texttt{envelope}. A reader rebuilds
  any shell as $|\mathbf q| = a\,g(c)$ (\cref{sec:shells}); a missing
  \texttt{egg} or \texttt{point} means 0. The attachment and top angles are
  not written.
\item \texttt{grid}: samples per petal, \texttt{nu} from base to tip and
  \texttt{nv} from one edge to the other. The core evaluates each sheet with
  2 samples per span in both directions, giving $18\cdot2+1 = 37$ by
  $10\cdot2+1 = 21$ points.
\item \texttt{channels}: the declaration of every latent channel (name, scope,
  dimensions, optional range, kind and, for categorical channels, labels;
  \cref{sec:channels}).
\item \texttt{petals}: id, layer and attachment angle of each petal.
\item \texttt{frames}: for each time $t$, a stage label, every petal's
  evaluated mid-surface points and thickness (row-major), and channel values.
  The closed bud is the frame at $t=0$. With development frames recorded, the
  growth stages come first, from $t=-1$ (first stage) up to $t=0$, each with
  its own \texttt{sepal\_radius}. Every channel's value is written in the first
  frame (per petal for petal channels, per frame for flower channels); a later
  frame that omits a channel keeps its last value. Every bud frame carries
  \texttt{stage: "bud"}.

  When the flower is opened (\cref{sec:flowering}), one frame per opening
  step is appended with \texttt{stage: "open"} and $t = k/N$ for
  $k = 1,\dots,N$ ($N = 24$ by default). Each petal entry holds its
  \texttt{mid} points, \texttt{thick}ness and \texttt{channels}, and each frame
  its flower-scope \texttt{channels}. The opening's exported channels
  (\texttt{vigour} per flower; \texttt{onset}, \texttt{rate},
  \texttt{length} and \texttt{width} per petal) are added to the top-level
  declaration, skipping names the bud already declared. The topology never
  changes and the across direction is fixed once on the bud, so the whole
  sequence plays back smoothly and exports as an animated glTF.
\item \texttt{meta}: the species (without its opening shape, which is not
  written), the plant's traits, the flower's traits and the development
  report.
\end{itemize}
All numbers are rounded to $10^{-5}$, far below petal
thickness, which keeps files small.

\subsection{Producing it}
The command-line tool writes one flower:
\begin{center}\small
\texttt{automark --seed N [--flower K] --json out.json [--species file.json]}\\
\texttt{[--relax N] [--preview] [--frames] [--open N]}
\end{center}
\texttt{--frames} records the development stages, \texttt{--open N} runs the
opening over $N$ frames, and \texttt{--species} reads a species JSON instead of
the built-in peony. \texttt{--preview} uses the fast preset (12 relaxation and
30 settling iterations instead of 30 and 75): about twice as fast, but its
geometry and channels differ slightly, so it is not canonical and not covered
by the golden digests (\cref{sec:determinism}). The browser build exposes the
same outputs (canonical bud, preview bud and opened flower).

\subsection{Orientation}
The surface normal $\partial_u \x \times \partial_v \x$ must point away from the
bud centre. For each petal the core sums
$(\partial_u\x\times\partial_v\x)\cdot\nabla\rt$ over its grid cells and, if the
sum is negative, reverses the across direction.

\subsection{Solid petals}
The viewer turns each mid-surface into a solid by offsetting it by $\pm t/2$
along the unit normal and closing the rim, so exported models have real
thickness.
